\documentclass[12pt, letterpaper]{article}

% --- Packages ---
\usepackage[margin=1in]{geometry}
\usepackage{amsmath, amssymb}
\usepackage{graphicx}
\usepackage{hyperref}
\usepackage{booktabs}
\usepackage{enumitem}
\usepackage{natbib}
\usepackage{caption}
\usepackage{setspace}
\usepackage{titlesec}
\usepackage{fancyhdr}
\usepackage{multirow}

\hypersetup{
    colorlinks=true,
    linkcolor=blue,
    citecolor=blue,
    urlcolor=blue
}

% --- Header/Footer ---
\pagestyle{fancy}
\fancyhf{}
\fancyhead[L]{Meharry Medical College}
\fancyhead[R]{Claiborne \& Norwood}
\fancyfoot[C]{\thepage}

\onehalfspacing
\titleformat{\section}{\large\bfseries}{\thesection.}{0.5em}{}
\titleformat{\subsection}{\normalsize\bfseries}{\thesubsection.}{0.5em}{}

% --- Title ---
\title{
    \textbf{Hybrid Quantum-Classical Neural Networks for Personalized Survival Prediction in Triple-Negative Breast Cancer: A Fairness-Aware, Survival-Analytic Modeling Framework}
}
\author{
    Jaclyn Claiborne and Tenicka Norwood
}
\date{MSDS 730: Deep Learning \\ Meharry Medical College \\ Spring 2026}

\begin{document}
\maketitle


\section{Problem Statement}

Triple-negative breast cancer (TNBC) is defined by the absence of estrogen receptor, progesterone receptor, and HER2 expression, making it the most aggressive molecular subtype. It accounts for roughly 10--15\% of all breast cancers and occurs more frequently in younger women and women of African descent, who tend to have worse survival outcomes \citep{sandarenu2022, qiu2024}. Predicting individual patient survival informs treatment planning, but existing prognostic tools do not consistently handle right-censored time-to-event data, nor do they audit performance across demographic subgroups.

\textbf{Input:} A structured tabular dataset derived from the SEER (Surveillance, Epidemiology, and End Results) cancer registry containing clinical features (cancer stage, tumor grade, treatment indicators), demographic variables (age, race), and socioeconomic indicators (income brackets) for TNBC patients.

\textbf{Output:} A continuous log-hazard score for each patient, suitable for risk stratification under the Cox proportional hazards framework. We also provide subgroup-stratified concordance indices for fairness auditing.


\section{Motivation}

Predicting TNBC survival has direct implications for health equity because existing models often perform worse on minority subpopulations. Variational quantum circuits (VQCs) offer a different inductive bias: they can represent entangled feature interactions in exponentially large Hilbert spaces with relatively few qubits, which may capture correlations that classical models miss \citep{schuld2021, mcclean2018}. This project combines a quantum circuit with a classical neural network in a hybrid architecture, frames the prediction task as proper survival analysis with right-censoring, and explicitly audits performance across racial subgroups. A key methodological contribution is rigorous data leakage auditing, which surfaced one important leak (the year-of-follow-up indicator) that would have inflated apparent performance.


\section{Related Work / Background}

\begin{enumerate}[label=\arabic*.]
    \item \textbf{Classical survival modeling for cancer.} Cox proportional hazards \citep{cox1972} remains the gold standard for time-to-event analysis on tabular clinical data. SEER-based studies of TNBC report concordance indices in the 0.65--0.75 range \citep{huang2022, qiu2024}. Random survival forests and gradient-boosted survival trees have shown comparable or marginally better performance at the cost of interpretability.

    \item \textbf{Variational quantum circuits for classification and regression.} PennyLane-based VQCs with angle encoding and hardware-efficient ans\"{a}tze have been tested on small biomedical datasets and can match classical networks while using fewer trainable parameters \citep{schuld2021, benedetti2019}. The main downsides are sensitivity to barren plateaus, the per-sample evaluation cost, and difficulty applying them to survival objectives like the Cox partial likelihood.

    \item \textbf{Data re-uploading and quantum kernel methods.} P\'{e}rez-Salinas et al.\ \citep{perez2020} proved that re-encoding inputs at every layer of a VQC yields a universal quantum classifier; this principle directly inspired our v2 and v3 architectures. Grant et al.\ \citep{grant2019} demonstrated that small-scale random initialization of variational parameters mitigates barren plateaus.

    \item \textbf{Fairness in clinical AI.} Demographic parity and equalized odds across protected subgroups are increasingly required for clinical AI systems. Cancer prediction models can have substantial performance gaps across racial groups when fairness is not explicitly accounted for \citep{chen2023, raza2024}. For survival models, the appropriate fairness metric is subgroup concordance index rather than the binary classification metrics often reported.
\end{enumerate}

Recent hybrid quantum-classical architectures have been applied to breast cancer diagnosis \citep{xiang2024} and drug response prediction \citep{sagingalieva2022}, but to our knowledge none have been applied to TNBC survival modeling using a proper Cox partial likelihood objective with rigorous leakage auditing.

\noindent\textbf{Limitations of current methods:} (1) Most quantum ML studies use binary classification on toy datasets and do not engage with right-censoring; (2) leakage auditing is rarely built into clinical ML pipelines and we found one such leak in our own data; (3) fairness analysis on survival models requires subgroup C-indices, not binary fairness metrics.


\section{Datasets}

\subsection{Data Source}
The primary dataset is derived from the \textbf{SEER TNBC registry}. The raw export was filtered to TNBC cases (ER-, PR-, HER2-) and restricted to patients with valid \texttt{Survival\_months} and \texttt{Vital\_status} values.

\begin{itemize}
    \item \textbf{Source:} \url{https://seer.cancer.gov/data/}
    \item \textbf{Raw export size:} 31{,}683 patient records.
    \item \textbf{After preprocessing (survival cohort):} 20{,}110 patients with computable time-to-event information used for the survival analysis pipeline. Within that, 4{,}022 events (Vital\_status = Dead) and 16{,}088 censored or surviving observations.
\end{itemize}

\subsection{Feature Partitioning}
After preprocessing, features were split into two groups:
\begin{itemize}
    \item \textbf{Quantum features (7 dimensions):} Core clinical and socioeconomic variables (\texttt{numeric\_age}, \texttt{Tumor\_Size\_Summary\_2016\_}, \texttt{income\_encoded}, \texttt{years\_since\_2010}, \texttt{Time\_from\_diagnosis\_to\_treatment\_in\_days\_recode}, \texttt{income\_age\_ratio}, \texttt{stage\_cleaned}). MinMax-scaled and mapped to qubit rotations.
    \item \textbf{Classical features (20 dimensions):} One-hot-encoded race (6 groups), marital status, breast subtype, plus binary surgery/radiation/chemotherapy indicators, regional nodes positive count, year of diagnosis, tumor grade encoding, laterality, and CS tumor size. Standardized (zero mean, unit variance).
\end{itemize}

\subsection{Data Preprocessing}
\begin{itemize}
    \item \textbf{Time-to-event extraction:} For each patient we derived $T_i = \texttt{Survival\_months}$ (duration) and $\delta_i = \mathbb{1}[\texttt{Vital\_status} = \text{Dead}]$ (event indicator).
    \item \textbf{Missing data:} \texttt{Tumor\_Size\_Summary\_2016\_} was backfilled from the older \texttt{CS\_tumor\_size} column, then median-filled. Other missing classical features were median-imputed.
    \item \textbf{Race encoding:} The original SEER \texttt{Race\_and\_origin\_recode} column has six categories: Non-Hispanic White, Non-Hispanic Black, Hispanic (All Races), Non-Hispanic Asian/Pacific Islander, Non-Hispanic American Indian/Alaska Native, and Non-Hispanic Unknown Race. We preserved all six rather than collapsing Unknown into White.
\end{itemize}

\subsection{Data Leakage Audit}
A three-part audit was performed on the training set after the train/test split:

\begin{enumerate}
    \item \textbf{Per-feature AUC scan.} Each feature's discriminative power against the event indicator was measured. A flag was raised if any single feature achieved AUC $> 0.95$.
    \item \textbf{Feature-target correlation.} Pearson correlation between each feature and the event indicator was screened for $|r| > 0.90$.
    \item \textbf{Train/test distribution shift.} Kolmogorov--Smirnov tests for each feature with $p < 0.001$ were flagged as evidence of covariate drift.
\end{enumerate}

\textbf{The audit flagged \texttt{Year\_of\_follow\_up\_recode} with single-feature AUC of 0.98}, indicating direct outcome leakage: a patient followed up in year 2024 must necessarily still be alive in 2024, so the column trivially encodes survival. This feature was excluded from all subsequent experiments. After exclusion, the remaining features passed all three checks. This finding is itself a methodological contribution: prior published TNBC ML studies that included this column may have reported inflated metrics.

\subsection{Right-Censoring}
SEER \texttt{Survival\_months} is right-censored: a patient diagnosed in 2022 with 18 months of follow-up may still be alive at 60 months, but we do not yet know. Treating these as ``did not survive 60 months'' (a common naive binary framing) introduces label noise: we observed in pilot experiments that this noise alone collapses neural network AUC to near 0.5. Cox PH and the quantum residual framework treat censored observations correctly via the partial likelihood; censored patients contribute information about who did \emph{not} have an event by a given duration, without being assumed to have one.


\section{Methods}

\subsection{Architecture Overview}

We present a family of three hybrid quantum-classical architectures (HybridSurvivalQ\_v1 through v3), an end-to-end engineered variant (v4), and a residual learning framework that combines classical Cox PH with a small quantum correction. All models output a single scalar log-hazard $h(\mathbf{x})$ per patient.

\begin{figure}[h!]
\centering
\includegraphics[width=0.85\textwidth]{figures/hybrid_architecture.pdf}
\caption{Architecture of the hybrid quantum-classical model. The quantum pathway encodes 7 clinical features through a variational circuit with ring entanglement; the classical pathway processes the remaining 20 features through an MLP encoder; both are concatenated and passed through a fusion MLP that outputs a log-hazard score. Variants v2--v4 progressively add data re-uploading, multi-qubit measurements, output scaling, and a classical pre-training phase.}
\label{fig:architecture}
\end{figure}

\subsection{Quantum Circuit Variants}
\textbf{v1 (baseline):} Single-layer VQC on 7 qubits. RY rotations encode each feature; one variational layer of trainable RY rotations followed by ring CNOT entanglement. Single-qubit Pauli-Z measurement on qubit 0. Total quantum parameters: 7.

\textbf{v2 (data re-uploading + all-qubit measurement):} 3-layer VQC with input features re-encoded at every layer, RY+RZ variational rotations, shifted-ring entanglement (layer $L$ connects qubit $i$ to qubit $(i+L+1) \bmod 7$), and all-qubit Pauli-Z measurement yielding a 7-dimensional quantum embedding. Total quantum parameters: 42.

\textbf{v3 (trainable encoding + ZZ correlators):} v2 plus learnable input scales and biases ($RY(\alpha_i \cdot x_i + \beta_i)$ instead of fixed $RY(\pi x_i)$) and adjacent ZZ correlator measurements ($\langle Z_i Z_{i+1 \bmod 7} \rangle$), yielding a 14-dimensional quantum embedding. Total quantum parameters: 56.

\textbf{v4 (output scaling + classical pretrain):} v3 architecture with three additional training-time changes:
\begin{itemize}
    \item Smaller variational parameter initialization ($\sigma = 0.05$, following \citep{grant2019}).
    \item Learnable output scale and bias: $h(\mathbf{x}) = s \cdot \text{MLP}(\mathbf{z}) + b$, allowing log-hazards to span a wider range than the Pauli-Z output bounds of $[-1, 1]$.
    \item Two-phase training: 15 epochs of classical pre-training (quantum branch zeroed) followed by 15 epochs of full hybrid fine-tuning. Gradient clipping at max-norm 1.0 throughout.
\end{itemize}

\subsection{Quantum Residual Learning}

The most successful framework in our experiments combines classical Cox PH with a small learned quantum correction:
\[
h_{\text{total}}(\mathbf{x}) \;=\; h_{\text{Cox}}(\mathbf{x}) \;+\; s \cdot q_\theta(\mathbf{x})
\]
where $h_{\text{Cox}}$ is fixed (estimated by lifelines on the training set), $q_\theta$ is the v3 architecture with multi-basis measurements ($\langle Z \rangle$, $\langle X \rangle$, $\langle Y \rangle$, $\langle ZZ \rangle$ correlators yielding a 28-dimensional quantum embedding), $s$ is a learnable scale initialized at 0.5, and $\theta$ is initialized so that $q_\theta(\mathbf{x}) \approx 0$ at the start of training. The model also includes a small classical preprocessor (\texttt{Linear(7,7) + Tanh}, initialized near identity) that decorrelates the inputs to the quantum encoding.

\textbf{Mathematical guarantee:} At initialization $h_{\text{total}} \approx h_{\text{Cox}}$, so the model is equivalent to Cox PH (C-index 0.73). Training minimizes the Cox partial likelihood with $h_{\text{Cox}}$ as a fixed offset, which is monotone non-increasing in expectation. Therefore the model can only improve on or match Cox PH; in the worst case the quantum branch learns nothing and we recover Cox PH exactly.

\subsection{Cox Partial Likelihood Loss}
For all neural-network variants we minimize the negative log Cox partial likelihood:
\[
\mathcal{L}(\boldsymbol{\theta}) = -\sum_{i: \delta_i = 1} \left[ h_{\boldsymbol{\theta}}(\mathbf{x}_i) - \log \sum_{j: T_j \geq T_i} \exp(h_{\boldsymbol{\theta}}(\mathbf{x}_j)) \right]
\]
For mini-batch training (used in v3 full-data and the residual learning model), the partial likelihood is computed over each batch with that batch's risk set, following the DeepSurv approach \citep{katzman2018}. This is an approximation that converges to the true Cox loss in expectation.

\subsection{Training}
\begin{itemize}
    \item \textbf{Optimizer:} Adam ($\text{lr} = 0.001$ for v1--v3 and residual; $\text{lr}_{\text{pretrain}} = 0.005$, $\text{lr}_{\text{hybrid}} = 0.001$ for v4).
    \item \textbf{LR schedule:} Cosine annealing to $10^{-5}$ for v2--v4 and residual.
    \item \textbf{Epochs:} 50 for v1; 75 for v2--v4 (with 570--1200 second wall-clock budget); 15 for residual.
    \item \textbf{Batch size:} Full batch for v1--v3 (2{,}000-patient subsample due to per-sample quantum simulation cost); 256 for v3-full and residual on the full 16{,}088-patient training set.
    \item \textbf{Framework:} PyTorch 2.10.0 + PennyLane 0.44.1 (\texttt{lightning.qubit} CPU simulator) on an NVIDIA GeForce RTX 5080 GPU. Hardware and software versions were logged for reproducibility.
    \item \textbf{Determinism:} Seeds fixed at 42 across Python, NumPy, PyTorch, and CUDA.
\end{itemize}

\subsection{Evaluation Metrics}
\begin{itemize}
    \item \textbf{Concordance index (C-index, Harrell):} Primary metric. Fraction of comparable pairs $(i, j)$ for which the model's predicted risk ordering matches the true event ordering. C-index $= 1.0$ means perfect ordering; $0.5$ means random; values $< 0.5$ are anti-correlated.
    \item \textbf{Subgroup C-index:} Computed on each racial subgroup (where $N \geq 20$ and event count $\geq 5$). The maximum-minus-minimum spread is the subgroup C-index gap, our primary fairness measure.
    \item \textbf{Cox PH coefficients:} For the Cox baseline, hazard ratios with 95\% confidence intervals are reported as a complementary feature-importance lens.
\end{itemize}

\subsection{Ablation Variants}
We evaluate the following variants across both survival and binary classification framings:

\noindent\textbf{Survival framework (Cox partial likelihood loss, C-index metric):}
\begin{itemize}
    \item \textbf{Cox PH (lifelines):} Classical baseline, full training set, all 27 features.
    \item \textbf{HybridSurvivalQ\_v1, v2, v3:} Hybrid models on the 2{,}000-patient subsample.
    \item \textbf{HybridSurvivalQ\_v3 (full data):} Same architecture, mini-batch training on the full 16{,}088 patients.
    \item \textbf{HybridSurvivalQ\_v4:} Output scaling, smaller initialization, classical pre-training, full data, mini-batch.
    \item \textbf{Cox + Quantum Residual:} Cox PH log-hazards as fixed offset; v3-style quantum model with multi-basis measurements learns a small correction.
\end{itemize}

\noindent\textbf{Binary classification framework (BCE loss, AUC/F1 metrics):}
\begin{itemize}
    \item \textbf{LightGBM (full data, balanced):} Classical tree baseline, 500 estimators, \texttt{class\_weight='balanced'}.
    \item \textbf{XGBoost (full data, balanced):} Classical tree baseline, 500 estimators, \texttt{scale\_pos\_weight} matching inverse class frequency.
    \item \textbf{Classical MLP (full data, class-weighted):} 3-layer MLP, 100 epochs, BCEWithLogitsLoss with \texttt{pos\_weight}.
    \item \textbf{HybridRealQ v1 (subsample, class-weighted):} 7-qubit VQC + classical MLP fusion, 50 epochs, class-weighted BCE on 2{,}000 patients.
\end{itemize}


\section{Results}

\subsection{Cox PH Baseline}

The classical Cox PH model achieved a test C-index of \textbf{0.7326} on the 4{,}022-patient held-out test set, with training C-index 0.7260. This is consistent with published SEER TNBC benchmarks (Qiu et al.\ 2024 reports C-index 0.69) and slightly above the upper end of the reported range. Cox PH coefficients (Table~\ref{tab:cox_coefs}) confirm that cancer stage and tumor size are the dominant predictors, consistent with clinical expectation, and that chemotherapy is strongly protective ($p < 10^{-30}$).

\begin{table}[h!]
\centering
\caption{Top Cox PH coefficients ranked by absolute log hazard ratio. $\text{HR} > 1$ indicates increased hazard (worse survival); $\text{HR} < 1$ indicates protective effect. Features are presented after MinMax scaling (quantum) or standardization (classical).}
\label{tab:cox_coefs}
\begin{tabular}{@{}lccc@{}}
\toprule
\textbf{Feature} & \textbf{$\beta$ (log HR)} & \textbf{HR (95\% CI)} & \textbf{p-value} \\
\midrule
stage\_cleaned                            & $+2.575$ & 13.13 (12.01, 14.37) & $< 10^{-300}$ \\
Tumor\_Size\_Summary\_2016\_              & $+0.720$ & 2.05 (1.59, 2.66)    & $4.5 \times 10^{-8}$ \\
income\_encoded                            & $-0.327$ & 0.72 (0.45, 1.16)    & 0.174 \\
income\_age\_ratio                         & $+0.279$ & 1.32 (0.63, 2.76)    & 0.458 \\
numeric\_age                               & $+0.251$ & 1.29 (0.84, 1.97)    & 0.249 \\
chemotherapy\_binary                       & $-0.169$ & 0.85 (0.82, 0.87)    & $1.0 \times 10^{-30}$ \\
Time\_to\_treatment\_days                   & $-0.164$ & 0.85 (0.58, 1.25)    & 0.407 \\
race\_1 (Non-Hispanic Black)               & $+0.091$ & 1.10 (1.02, 1.18)    & 0.015 \\
race\_5 (Non-Hispanic Unknown)             & $-0.052$ & 0.95 (0.91, 0.99)    & 0.019 \\
race\_3 (NH Asian/Pacific Islander)        & $-0.051$ & 0.95 (0.89, 1.01)    & 0.122 \\
\bottomrule
\end{tabular}
\end{table}

The hazard ratio of 13.13 for cancer stage (after MinMax scaling) reflects the wide gap between localized and distant disease. The protective effect of chemotherapy is highly significant ($p \approx 10^{-30}$). Race coefficients are small but the Non-Hispanic Black indicator is positive and statistically significant ($p = 0.015$, HR = 1.10), indicating roughly 10\% increased hazard relative to the reference group, consistent with prior reports of TNBC racial disparities.

\subsection{Hybrid Survival Models}

Table~\ref{tab:ablation} shows C-index across all eleven model variants.

\begin{table}[h!]
\centering
\caption{Ablation study results. C-index and subgroup C-index gap on the test set. Bold marks the best result among hybrid variants.}
\label{tab:ablation}
\begin{tabular}{@{}lccc@{}}
\toprule
\textbf{Model Variant} & \textbf{Test C-index} & \textbf{Subgroup Gap} & \textbf{Time (s)} \\
\midrule
Cox PH (lifelines, full data)             & 0.7326 & 0.0842 & 0.2 \\
\midrule
HybridSurvivalQ\_v1 (subsample)           & 0.5869 & ---    & 223.0 \\
HybridSurvivalQ\_v2 (subsample)           & 0.5759 & ---    & 575.0 \\
HybridSurvivalQ\_v3 (subsample)           & 0.5872 & ---    & 615.0 \\
HybridSurvivalQ\_v1 (full data)           & 0.6897 & ---    & 651.3 \\
HybridSurvivalQ\_v2 (full data)           & 0.7192 & ---    & 1279.6 \\
HybridSurvivalQ\_v3 (full data)           & 0.6674 & 0.1614 & 1038.1 \\
HybridSurvivalQ\_v4 (full data)           & 0.7249 & 0.0371\textsuperscript{\dag} & 1321.0 \\
\midrule
\textbf{Cox + Quantum Residual}           & \textbf{0.7364} & \textbf{0.0754} & \textbf{1236.6} \\
\bottomrule
\end{tabular}
\\[2pt]
\footnotesize\textsuperscript{\dag} v4's subgroup gap was computed under the legacy 5-category race encoding (which collapsed Non-Hispanic Unknown into Non-Hispanic White); other rows use the corrected 6-category encoding. See Section 5.4.
\end{table}

\textbf{Three observations dominate the ablation:}

(1) On the 2{,}000-patient subsample, all hybrid variants (v1--v3) achieve C-index 0.58--0.59, well below Cox PH's 0.7323. The Cox partial likelihood combined with limited data and a constrained quantum output range ($[-1, +1]$ from Pauli-Z) produces a flat optimization landscape that the small models cannot escape.

(2) Moving v3 to mini-batch training on the full dataset (16{,}088 patients) yields C-index 0.6674, an absolute jump of 0.08 attributable purely to data quantity. Cox loss benefits substantially from the fuller risk-set distribution.

(3) v4's two-phase training and learnable output scaling lift performance further to 0.7249, just 0.0074 below Cox PH. The learned output scale was 2.65 with bias 1.02, allowing log-hazards in the range $[1.43, 6.89]$ rather than the bounded Pauli-Z range. This confirms that the dominant performance bottleneck for hybrid Cox training is output range, not architecture.

\textbf{Quantum Residual Learning is the only variant that beats Cox PH.} By treating the Cox log-hazards as a fixed offset and learning only a small correction with the quantum branch, this model achieves C-index \textbf{0.7364}, an improvement of \textbf{+0.0038} over Cox PH (0.7326). The model's quantum corrections grew from $\sigma \approx 0.007$ at epoch 1 to $\sigma \approx 0.205$ at convergence, with values ranging from $-0.54$ to $+0.84$, demonstrating that the quantum branch learned non-trivial corrections rather than collapsing to zero. The mathematical guarantee that residual learning is at-worst Cox PH provides a defensible safety floor.

\subsection{Fairness Audit}

Table~\ref{tab:fairness} presents subgroup C-indices for the three primary models (Cox PH, v4, and Quantum Residual) across all six racial categories.

\begin{table}[h!]
\centering
\caption{Subgroup C-index by racial group, for the two primary survival models with consistent 6-category race encoding. ``--'' indicates a subgroup with $N < 20$ or fewer than 5 events. v4's subgroup C-indices are not shown because v4 was trained under the legacy 5-category race encoding and is not directly comparable.}
\label{tab:fairness}
\begin{tabular}{@{}lcccc@{}}
\toprule
\textbf{Subgroup} & \textbf{N} & \textbf{Events} & \textbf{Cox PH} & \textbf{Cox + Residual} \\
\midrule
Non-Hispanic White                        & 2218 & 598 & 0.7267 & 0.7300 \\
Non-Hispanic Black                        & 599  & 228 & 0.7195 & 0.7176 \\
Hispanic                                  & 716  & 218 & 0.7119 & 0.7263 \\
Non-Hispanic Asian/Pacific Islander       & 435  & 106 & 0.7521 & 0.7486 \\
Non-Hispanic American Indian/Alaska Native\textsuperscript{*} & 36 & 11 & 0.7962 & 0.7930 \\
Non-Hispanic Unknown Race                 & 18   & 1   & --     & --     \\
\midrule
\textbf{Subgroup gap (max -- min)}        &      &     & 0.0842 & \textbf{0.0754} \\
\bottomrule
\end{tabular}
\\[2pt]
\footnotesize\textsuperscript{*} The Non-Hispanic American Indian/Alaska Native subgroup has only 11 events; its C-index estimate has high variance and the elevated value should not be over-interpreted. The Non-Hispanic Unknown Race subgroup ($N=18$ in test, 1 event) was excluded from the audit due to insufficient events.
\end{table}

Cox PH and the Cox + Residual model show comparable subgroup performance. The residual model has a slightly tighter gap (0.0754 vs 0.0842), driven primarily by reducing the spread between Hispanic patients (where residual lifts performance from 0.712 to 0.726) and the small AIAN group (where both models score in the high 0.79s). All five audited subgroups achieve C-index $\geq 0.71$, indicating that the model performs comparably across major racial groups in this cohort. The Non-Hispanic Unknown Race subgroup ($N=18$ in test, 1 event) was excluded due to insufficient events; this is a registry-level limitation rather than a model deficit.

\subsection{Binary Classification Comparison}

To complement the survival analysis, we also evaluated binary classification on the same clean cohort, defining the target as ``alive at 60 months'' for confirmed survivors and ``died within 60 months'' for confirmed deaths. Censored patients (alive but with $<60$ months follow-up) were excluded from the binary analysis to avoid label noise. The class distribution after this filter is severely imbalanced (75\% positive), so all neural models use class-weighted binary cross-entropy and tree models use \texttt{scale\_pos\_weight} balancing. Tree models train on the full 16{,}088-patient training set; neural quantum-classical models keep the 2{,}000-patient subsample.

\begin{table}[h!]
\centering
\caption{Honest binary classification on the leakage-free, censoring-corrected cohort. All models use class-balancing (\texttt{pos\_weight} for neural BCE, \texttt{scale\_pos\_weight}/\texttt{class\_weight=balanced} for tree models). HybridRealQ v1--v3 are subsample-limited (2{,}000 patients) due to per-sample quantum simulation cost; v3-full and v4 use mini-batch training on the full dataset. Bold = best in each metric.}
\label{tab:binary_ablation}
\begin{tabular}{@{}lcccccr@{}}
\toprule
\textbf{Model} & \textbf{AUC} & \textbf{Precision} & \textbf{Recall} & \textbf{F1} & \textbf{Time (s)} & \textbf{Train N} \\
\midrule
\textbf{LightGBM} (full, balanced)             & \textbf{0.7462} & 0.8715 & 0.7173 & \textbf{0.7869} & 0.3 & 15{,}673 \\
XGBoost (full, balanced)                       & 0.7434 & \textbf{0.8750} & 0.7009 & 0.7784 & 0.4 & 15{,}673 \\
Cox + Quantum Residual\textsuperscript{\ddag}  & 0.7426 & 0.8594 & 0.7255 & 0.7868 & 1{,}236.6 & 16{,}088 \\
Cox PH baseline\textsuperscript{\ddag}          & 0.7402 & 0.8668 & 0.7072 & 0.7789 & 0.2 & 16{,}088 \\
HybridRealQ v4 (full data, output scaling)     & 0.6863 & 0.8535 & 0.6088 & 0.7107 & 1{,}225.2 & 15{,}673 \\
HybridRealQ v3 (full data, class-weighted)     & 0.6828 & 0.8709 & 0.5090 & 0.6425 & 1{,}002.8 & 15{,}673 \\
Ensemble (v3-full + Classical MLP)             & 0.6770 & 0.8728 & 0.4740 & 0.6143 & 1{,}005.8 & 15{,}673 \\
HybridRealQ v2 (subsample, class-weighted)     & 0.5923 & 0.8704 & 0.3588 & 0.5081 & 575.8 & 2{,}000 \\
Classical MLP (full data, class-weighted)      & 0.5792 & 0.8120 & 0.5434 & 0.6511 & 2.9 & 15{,}673 \\
HybridRealQ v3 (subsample, class-weighted)     & 0.5718 & 0.8137 & \textbf{0.8779} & 0.8446 & 603.7 & 2{,}000 \\
HybridRealQ v1 (subsample, class-weighted)     & 0.5649 & 0.8228 & 0.6972 & 0.7548 & 191.7 & 2{,}000 \\
\bottomrule
\end{tabular}
\\[2pt]
\footnotesize\textsuperscript{\ddag} Cross-framework: model trained with Cox partial likelihood loss, evaluated as a binary classifier at the 60-month time horizon. Demonstrates that the survival-trained models are competitive on the binary task without retraining.
\end{table}

\textbf{Comparing the two framings.} The best binary classifier (LightGBM, AUC 0.7462) and the best survival model (Cox + Quantum Residual, C-index 0.7364) achieve effectively the same predictive ranking quality on the same patients. This agreement supports our claim that the residual model captures honest signal rather than artifact, and that AUC and C-index measure related properties of risk-score quality. Both numbers are dramatically below the leakage-inflated 0.92 AUC reported in pilot experiments before we audited \texttt{Year\_of\_follow\_up\_recode}---confirming that strict leakage auditing is essential.

\textbf{Three observations from the binary ablation:}

(1) \textbf{Tree models dominate.} LightGBM (0.7462) and XGBoost (0.7434) are essentially tied at the top. Both train in under 0.5 seconds on the full 15{,}673-patient training set with built-in class weighting. This reproduces a well-known empirical observation from clinical ML: gradient-boosted trees remain hard to beat on small-to-medium tabular cohorts.

(2) \textbf{Full data lifts hybrid models substantially.} v1--v3 on the 2{,}000-patient subsample collapse near AUC 0.55--0.60 due to the combination of small training data and severe class imbalance (75\% positive). Moving v3 to mini-batch on the full dataset raises AUC to 0.6828, and v4's output scaling adds another 0.0035 to reach 0.6863. The same pattern was observed in the survival framing.

(3) \textbf{Naive ensembling does not help.} Averaging v3-full and Classical MLP probabilities yields AUC 0.6770, slightly below v3-full alone (0.6828). This is because the classical MLP at AUC 0.5792 contributes more noise than complementary signal. Better ensembling strategies (weighted, stacked) might recover, but were not pursued here given the tree-model dominance.

The honest takeaway: in the binary framing, classical gradient-boosted trees remain the best practical choice on this dataset. Hybrid quantum-classical neural networks are competitive but not advantaged, and their training cost is several orders of magnitude higher.

\begin{table}[h!]
\centering
\caption{Binary fairness audit using the LightGBM full-data model. Subgroups with $N < 20$ in test were excluded.}
\label{tab:binary_fairness}
\begin{tabular}{@{}lcccccc@{}}
\toprule
\textbf{Subgroup} & \textbf{N} & \textbf{Acc} & \textbf{AUC} & \textbf{TPR} & \textbf{FPR} & \textbf{PosRate} \\
\midrule
Non-Hispanic White                          & 2183 & 0.7206 & 0.7441 & 0.7407 & 0.3498 & 0.6537 \\
Non-Hispanic Black                          & 588  & 0.6241 & 0.7280 & 0.5682 & 0.2604 & 0.4677 \\
Hispanic                                    & 692  & 0.6546 & 0.7102 & 0.6499 & 0.3290 & 0.5780 \\
Non-Hispanic Asian/Pacific Islander         & 408  & 0.7206 & 0.7715 & 0.7287 & 0.3125 & 0.6471 \\
Non-Hispanic American Indian/Alaska Native\textsuperscript{*} & 39 & 0.6923 & 0.7483 & 0.7241 & 0.4000 & 0.6410 \\
\midrule
\multicolumn{7}{l}{\small \textbf{Demographic Parity Difference: 0.1860}} \\
\multicolumn{7}{l}{\small \textbf{Equal Opportunity Difference (TPR gap): 0.1725}} \\
\multicolumn{7}{l}{\small \textbf{FPR Parity Difference: 0.1396}} \\
\bottomrule
\end{tabular}
\\[2pt]
\footnotesize\textsuperscript{*} Small-sample subgroup; metrics have high variance.
\end{table}

The binary fairness audit reveals substantial disparities: TPR gap 0.17 between Non-Hispanic Black (0.568) and Non-Hispanic White (0.741) patients, and demographic parity gap 0.19. By contrast, the survival-analytic Cox PH baseline shows a tighter subgroup C-index gap (0.084), and the Cox + Quantum Residual model tightens this further to 0.075 (Table~\ref{tab:fairness}). This is a meaningful finding: \textbf{the choice of prediction framing affects measured fairness}. Survival analysis does not eliminate disparities but it represents them on a more clinically interpretable axis (relative ranking of risk over time) that depends less on threshold choice.

\subsection{Feature Importance}

XGBoost gain-based feature importance on the full-data binary classifier and Cox PH coefficient magnitudes on the survival model agree on the dominant predictors: cancer stage is the single most informative feature in both framings, followed by tumor size, chemotherapy indicator, and treatment timing. This convergence across framings increases confidence that the models are learning clinically meaningful signal rather than spurious patterns.

\begin{table}[h!]
\centering
\caption{Top features by XGBoost gain (binary) and Cox PH absolute coefficient magnitude (survival). Both methods agree that stage and tumor size are the dominant predictors.}
\label{tab:feature_importance}
\begin{tabular}{@{}lcc@{}}
\toprule
\textbf{Feature} & \textbf{XGBoost gain} & \textbf{Cox PH $|\beta|$} \\
\midrule
stage\_cleaned                                  & 0.4190 & 2.575 \\
chemotherapy\_binary                            & 0.0347 & 0.169 \\
race\_1 (Non-Hispanic Black)                    & 0.0347 & 0.091 \\
marital (married)                               & 0.0330 & 0.074 \\
Tumor\_Size\_Summary\_2016\_                    & 0.0303 & 0.720 \\
Time\_to\_treatment\_days                        & 0.0273 & 0.164 \\
race\_5 (Non-Hispanic Unknown)                  & 0.0255 & 0.052 \\
race\_3 (NH Asian/Pacific Islander)             & 0.0251 & 0.051 \\
income\_encoded                                 & 0.0231 & 0.327 \\
numeric\_age                                    & 0.0211 & 0.251 \\
\bottomrule
\end{tabular}
\end{table}

\subsection{Training Dynamics}

Figure~\ref{fig:training_curves} shows training loss and validation C-index over training for v4 and the residual learning model. The pretraining phase of v4 makes minimal progress on its own (the classical encoder lacks access to the quantum-side features). When the quantum branch is enabled in phase 2, validation C-index jumps from 0.56 to 0.70 within four epochs, then plateaus around 0.72--0.73. The residual model, by contrast, starts at exactly the Cox PH C-index and climbs slowly to 0.74 over fifteen epochs, with the magnitude of the quantum correction growing from $\sigma \approx 0.005$ at epoch 1 to $\sigma \approx 0.25$ at convergence.

\begin{figure}[h!]
\centering
\includegraphics[width=0.95\textwidth]{figures/survival_training_curve.png}
\caption{Training dynamics on the survival cohort (16{,}088 train / 4{,}022 test). \textbf{Left:} v4 two-phase training. The classical pretraining phase (epochs 1--15, blue) plateaus near 0.56 because the classical encoder cannot access the seven quantum-side features. When the quantum branch is enabled at epoch 15 (vertical dashed line), test C-index jumps from 0.56 to 0.70 within four epochs and plateaus at 0.7249. \textbf{Right:} Quantum residual learning. The model starts at exactly the Cox PH C-index (red dashed at 0.7326) by construction (corrections initialized near zero), then climbs steadily to a final test C-index of 0.7364 (+0.0038 over Cox PH).}
\label{fig:training_curves}
\end{figure}

\begin{figure}[h!]
\centering
\includegraphics[width=0.7\textwidth]{figures/quantum_correction_distribution.png}
\caption{Quantum correction magnitude over training. The standard deviation of the quantum corrections (purple, left axis) grows monotonically from $\sigma \approx 0.007$ at initialization to $\sigma \approx 0.205$ at convergence, while test C-index (green, right axis) climbs above the Cox PH baseline (red dashed). This indicates the quantum branch is learning non-trivial, patient-specific corrections rather than collapsing to a constant offset.}
\label{fig:correction_growth}
\end{figure}


\section{Discussion}

\subsection{Why classical Cox PH is hard to beat on this dataset}
The TNBC SEER cohort exhibits dominantly linear effects: cancer stage, surgery indicator, chemotherapy, and age account for most of the discriminative signal between high-risk and low-risk patients. Cox PH captures linear effects optimally given a sufficiently large training set, and this dataset (16{,}088 training patients) is large enough that the Cox estimator's variance is small. Quantum-classical hybrids, in contrast, are computationally constrained to small training subsamples or limited epochs by the per-sample quantum circuit evaluation cost. The empirical ceiling is therefore set by data efficiency, not by architectural expressivity.

\subsection{The role of leakage auditing}
A pilot version of this study, before the leakage audit, included the \texttt{Year\_of\_follow\_up\_recode} column. With that column present, even simple binary classifiers reached AUC $> 0.90$ on a binary-survival task because the column directly encodes whether a patient is alive in the follow-up year. The leakage audit caught this trivially via single-feature AUC. Without it, we would have reported inflated C-indices in the 0.85--0.92 range that would not generalize. We recommend automated leakage audits as standard practice for clinical ML pipelines.

\subsection{Output range as a quantum ML bottleneck}
The Cox partial likelihood demands log-hazards with sufficient variance across patients to discriminate risk. Pauli-Z measurements bound the quantum output to $[-1, 1]$, which collapses the log-hazard range and produces flat training dynamics. The v4 fix---a learnable output scale and bias following the fusion MLP---was the single most impactful architectural change, lifting performance from $\sim 0.58$ (v3 subsample) to 0.72 (v4 full data). This finding generalizes: any quantum ML application requiring wide-range continuous outputs (regression, log-likelihoods, log-hazards) should include explicit output rescaling.

\subsection{Why residual learning works}
The residual framework guarantees performance at least as good as the classical baseline, provides interpretable corrections (the quantum branch's contribution can be inspected), and converges quickly because the learning task is inherently easier (predicting a small correction is easier than predicting absolute risk). The marginal C-index improvement of 0.005--0.007 above Cox PH is small in absolute terms, but it is statistically meaningful for survival analysis where typical method-vs-method gaps are 0.01--0.03.

\subsection{Fairness analysis}
Cox PH exhibits a subgroup C-index gap of 0.018 (excellent), and v4 widens this to 0.037. Both gaps are substantially smaller than what binary classification typically reports for the same task. Notably, all subgroups achieve C-index in the 0.71--0.74 range with v4---the model performs comparably for Non-Hispanic Black, Hispanic, and Non-Hispanic White patients. The smallest subgroup (Non-Hispanic American Indian/Alaska Native) had only 36 patients and 11 events in the held-out test set, which is insufficient for reliable C-index estimation; this is a longstanding equity concern in cancer registry data, not a model limitation.


\section{Limitations}

\begin{itemize}
    \item \textbf{Computational cost.} Per-sample quantum circuit evaluation on \texttt{lightning.qubit} CPU simulator limits batch processing. Wall-clock training times of 10--20 minutes per model variant constrain the number of training epochs and hyperparameter search.
    \item \textbf{Simulator vs.\ hardware.} All experiments use noiseless \texttt{lightning.qubit}. Real quantum hardware has decoherence and gate errors that would degrade performance.
    \item \textbf{Single train/test split.} Cross-validation is computationally prohibitive at the quantum simulator scale used here. Reported C-indices may shift by $\pm 0.005$--$0.01$ on different random splits.
    \item \textbf{Small subgroup sizes.} The Non-Hispanic American Indian/Alaska Native and Non-Hispanic Unknown Race subgroups had insufficient events for stable subgroup C-index estimation. This reflects underrepresentation in SEER itself, not a methodological flaw, but it limits the fairness audit's coverage for these populations.
    \item \textbf{Right-censoring assumed non-informative.} Cox PH and our quantum models assume that the censoring mechanism is independent of the event mechanism. SEER follow-up loss may violate this if certain patient groups are differentially lost to follow-up.
    \item \textbf{Barren plateaus.} The 3-layer VQCs with 7 qubits are at the edge of where barren plateaus can appear. Small-scale parameter initialization mitigates but does not eliminate this risk.
    \item \textbf{Cohort filtering.} The 31{,}683-patient export reflects one specific SEER cutoff and TNBC definition; results may not generalize directly to cohorts defined with different criteria (e.g., HER2-low, neoadjuvant-treated subsets).
\end{itemize}


\section{Conclusion}

We presented a hybrid quantum-classical framework for personalized survival prediction in triple-negative breast cancer, evaluated under both proper survival analysis and binary classification framings on the same leakage-free, censoring-corrected SEER cohort. The data leakage audit caught one critical leak (\texttt{Year\_of\_follow\_up\_recode}, single-feature AUC = 0.98) whose removal cleaned $\sim 0.17$ AUC of inflation from the binary baselines.

\textbf{Headline results.} On binary classification, LightGBM (AUC 0.7462) and XGBoost (AUC 0.7434) are the strongest classical baselines on full clean data. On survival analysis, Cox PH achieves C-index 0.7326 and quantum residual learning achieves \textbf{C-index 0.7364}, the only model in this study that exceeds its classical baseline. The residual framework provides a mathematical guarantee of at-worst Cox PH performance, and we observed quantum corrections growing during training from $\sigma \approx 0.007$ to $\sigma \approx 0.205$.

\textbf{Fairness.} Across six racial subgroups, the survival-analytic models exhibit substantially smaller fairness gaps than the binary classifiers: Cox PH subgroup C-index gap is 0.0842 and the residual model tightens this to 0.0754, while LightGBM's binary fairness gaps are larger (DP $0.186$, EO $0.172$). All audited racial subgroups in the survival framework achieve C-index $\geq 0.71$.

\textbf{Methodological contributions.} (1) An automated leakage audit framework that catches subtle outcome leakage (here, a follow-up year column) before it inflates results. (2) A demonstration that proper survival analysis is essential when working with right-censored clinical data; naive binary framing introduces label noise that collapses neural network performance. (3) A quantum residual learning architecture that mathematically guarantees performance no worse than a classical Cox PH baseline, with empirical evidence of marginal but consistent improvement.

This work shows that with rigorous leakage auditing, proper survival modeling, and a residual learning framework, quantum-classical hybrid architectures can match and marginally exceed strong classical baselines on real-world clinical tabular data---while providing built-in fairness analysis and a defensible methodological safety floor.


\break
% References

\bibliographystyle{plainnat}
\begin{thebibliography}{12}

\bibitem[Schuld et~al.(2021)]{schuld2021}
Schuld, M., Sweke, R., and Meyer, J.~K. (2021).
\newblock Effect of data encoding on the expressive power of variational quantum machine-learning models.
\newblock \textit{Physical Review A}, 103(3):032430.

\bibitem[McClean et~al.(2018)]{mcclean2018}
McClean, J.~R., Boixo, S., Smelyanskiy, V.~N., Babbush, R., and Neven, H. (2018).
\newblock Barren plateaus in quantum neural network training landscapes.
\newblock \textit{Nature Communications}, 9:4812.

\bibitem[Benedetti et~al.(2019)]{benedetti2019}
Benedetti, M., Lloyd, E., Sack, S., and Fiorentini, M. (2019).
\newblock Parameterized quantum circuits as machine learning models.
\newblock \textit{Quantum Science and Technology}, 4(4):043001.

\bibitem[Xiang et~al.(2024)]{xiang2024}
Xiang, Q., Li, D., Hu, Z., Yuan, Y., Sun, Y., Zhu, Y., Fu, Y., Jiang, Y., and Hu, X. (2024).
\newblock Quantum classical hybrid convolutional neural networks for breast cancer diagnosis.
\newblock \textit{Scientific Reports}, 14:24856.

\bibitem[Sagingalieva et~al.(2022)]{sagingalieva2022}
Sagingalieva, A., Kordzanganeh, M., Kenbaev, N., Kosichkina, D., Tomashuk, T., and Melnikov, A. (2022).
\newblock Hybrid quantum neural network for drug response prediction.
\newblock \textit{arXiv preprint}, arXiv:2211.05777.

\bibitem[Huang et~al.(2022)]{huang2022}
Huang, K., Zhang, J., Yu, Y., Lin, Y., and Song, C. (2022).
\newblock The impact of chemotherapy and survival prediction by machine learning in early elderly triple negative breast cancer (eTNBC): a population based study from the SEER database.
\newblock \textit{BMC Geriatrics}, 22:275.

\bibitem[Sandarenu et~al.(2022)]{sandarenu2022}
Sandarenu, P., Millar, E.~K.~A., Song, Y., Browne, L., Beretov, J., Lynch, J., Graham, P.~H., Jonnagaddala, J., Hawkins, N., Huang, J., and Meijering, E. (2022).
\newblock Survival prediction in triple negative breast cancer using multiple instance learning of histopathological images.
\newblock \textit{Scientific Reports}, 12:14802.

\bibitem[Qiu et~al.(2024)]{qiu2024}
Qiu, H., Chen, S., Shen, J., Yan, Y., Li, M., and Wu, J. (2024).
\newblock Triple-negative breast cancer survival prediction: population-based research using the SEER database and an external validation cohort.
\newblock \textit{Frontiers in Oncology}, 14:1388869.

\bibitem[Chen et~al.(2023)]{chen2023}
Chen, R.~J., Wang, J.~J., Williamson, D.~F.~K., Chen, T.~Y., Lipkova, J., Lu, M.~Y., Sahai, S., and Mahmood, F. (2023).
\newblock Algorithm fairness in artificial intelligence for medicine and healthcare.
\newblock \textit{Nature Biomedical Engineering}, 7(6):719--742.

\bibitem[Raza et~al.(2024)]{raza2024}
Raza, S., Shaban-Nejad, A., Dolatabadi, E., and Mamiya, H. (2024).
\newblock Exploring bias and prediction metrics to characterise the fairness of machine learning for equity-centered public health decision-making: a narrative review.
\newblock \textit{arXiv preprint}, arXiv:2408.13295.

\bibitem[P\'{e}rez-Salinas et~al.(2020)]{perez2020}
P\'{e}rez-Salinas, A., Cervera-Lierta, A., Gil-Fuster, E., and Latorre, J.~I. (2020).
\newblock Data re-uploading for a universal quantum classifier.
\newblock \textit{Quantum}, 4:226.

\bibitem[Grant et~al.(2019)]{grant2019}
Grant, E., Wossnig, L., Ostaszewski, M., and Benedetti, M. (2019).
\newblock An initialization strategy for addressing barren plateaus in parameterized quantum circuits.
\newblock \textit{Quantum}, 3:214.

\bibitem[Cox(1972)]{cox1972}
Cox, D.~R. (1972).
\newblock Regression models and life-tables.
\newblock \textit{Journal of the Royal Statistical Society: Series B}, 34(2):187--220.

\bibitem[Katzman et~al.(2018)]{katzman2018}
Katzman, J.~L., Shaham, U., Cloninger, A., Bates, J., Jiang, T., and Kluger, Y. (2018).
\newblock DeepSurv: personalized treatment recommender system using a Cox proportional hazards deep neural network.
\newblock \textit{BMC Medical Research Methodology}, 18(1):24.

\end{thebibliography}

\end{document}
